%% Chương 6: Kết quả và Thảo luận
\clearpage
\phantomsection
\section*{CHƯƠNG 6. KẾT QUẢ VÀ THẢO LUẬN}
\addcontentsline{toc}{section}{\numberline{}CHƯƠNG 6. KẾT QUẢ VÀ THẢO LUẬN}
\setcounter{section}{6}
\setcounter{subsection}{0}
\setcounter{figure}{0}
\setcounter{table}{0}
\setcounter{equation}{0}

%% Bảng metrics — trích từ summary_tables_testkey.md (key: testkey, 2026-06-23)
%% File này được \input từ ch6_results.tex (section counter = 6)
%% Bảng tự đánh số 6.1, 6.2, ... theo \thetable = \thesection.\arabic{table}

%% ============================================================
%% Bảng 1: Tổng quan
%% ============================================================
\begin{table}[ht]
\centering
\caption{Tổng quan thực nghiệm — file output.h264}
\label{tab:overview}
\begin{tabular}{lll}
\toprule
 & S1 (Type 1 + Type 5) & S2 (Type 5 only) \\
\midrule
NALUs mã hoá       & 19{,}037 & 155 \\
Thuật toán         & hybrid, aes, des, rc4 & hybrid, aes, des, rc4 \\
Byte-exact verify  & 8/8 PASS & 8/8 PASS \\
\bottomrule
\end{tabular}
\end{table}

%% ============================================================
%% Bảng 2: Performance S1
%% ============================================================
\begin{table}[ht]
\centering
\caption{Hiệu năng — Chiến lược S1 (19{,}037 NALUs, key: testkey)}
\label{tab:perf_s1}
\begin{tabular}{lrrrr}
\toprule
Metric & hybrid & aes & des & rc4 \\
\midrule
Encrypt time (pipeline, s)  & 1.510 & 0.984 & 1.519 & 1.082 \\
DC encrypt time (ms)        & 556.751 & 52.416 & 447.375 & 14.649 \\
DC encrypt / NALU (ns)      & 29{,}245 & 2{,}753 & 23{,}500 & 769 \\
Decrypt time (pipeline, s)  & 1.633 & 1.008 & 8.686$^\dagger$ & 0.949 \\
DC decrypt time (ms)        & 692.768 & 55.970 & 478.312 & 10.925 \\
DC decrypt / NALU (ns)      & 36{,}390 & 2{,}940 & 25{,}125 & 573 \\
\bottomrule
\end{tabular}
\begin{tablenotes}
\small
\item $^\dagger$ Pipeline anomaly hệ thống (disk I/O). DC decrypt thực tế = 478 ms, bình thường.
\end{tablenotes}
\end{table}

%% ============================================================
%% Bảng 3: Performance S2
%% ============================================================
\begin{table}[ht]
\centering
\caption{Hiệu năng — Chiến lược S2 (155 NALUs, key: testkey)}
\label{tab:perf_s2}
\begin{tabular}{lrrrr}
\toprule
Metric & hybrid & aes & des & rc4 \\
\midrule
Encrypt time (pipeline, s)  & 0.990 & 0.837 & 1.109 & 0.960 \\
DC encrypt time (ms)        & 6.659 & 0.477 & 3.475 & 0.112 \\
DC encrypt / NALU (ns)      & 42{,}962 & 3{,}079 & 22{,}417 & 724 \\
DC chiếm \% pipeline encrypt & 0.67\% & 0.057\% & 0.31\% & 0.012\% \\
Decrypt time (pipeline, s)  & 0.937 & 0.873 & 1.051 & 0.907 \\
DC decrypt time (ms)        & 7.215 & 0.475 & 3.640 & 0.081 \\
DC decrypt / NALU (ns)      & 46{,}547 & 3{,}063 & 23{,}484 & 525 \\
\bottomrule
\end{tabular}
\begin{tablenotes}
\small
\item Lưu ý: Pipeline time hội tụ (0.84--1.11 s) vì DC $< 1\%$ tổng thời gian;
I/O đọc/ghi 219 MB chiếm $> 99\%$.
\end{tablenotes}
\end{table}

%% ============================================================
%% Bảng tỉ lệ DC trong pipeline S1
%% ============================================================
\begin{table}[ht]
\centering
\caption{DC encrypt time chiếm \% pipeline — S1 vs S2}
\label{tab:dc_percent}
\begin{tabular}{lrr}
\toprule
Thuật toán & DC / pipeline S1 & DC / pipeline S2 \\
\midrule
hybrid & 36.9\% & 0.67\% \\
aes    &  5.3\% & 0.057\% \\
des    & 29.5\% & 0.31\% \\
rc4    &  1.35\% & 0.012\% \\
\bottomrule
\end{tabular}
\end{table}

%% ============================================================
%% Bảng 4: Security
%% ============================================================
\begin{table}[ht]
\centering
\caption{Chỉ số bảo mật — S1 và S2 (key: testkey)}
\label{tab:security}
\resizebox{\textwidth}{!}{%
\begin{tabular}{llrrrrrrrrr}
\toprule
 & & \multicolumn{4}{c}{S1 (19{,}037 NALUs)} & & \multicolumn{4}{c}{S2 (155 NALUs)} \\
\cmidrule{3-6} \cmidrule{8-11}
Metric & Ngưỡng & hybrid & aes & des & rc4 & & hybrid & aes & des & rc4 \\
\midrule
NPCR plaintext (\%) & $> 99.6$ & 99.5260 & N/A & N/A & N/A & & 99.6903 & N/A & N/A & N/A \\
UACI plaintext (\%) & $\approx 33.46$ & 32.7414 & N/A & N/A & N/A & & 32.4751 & N/A & N/A & N/A \\
Entropy DC (bits/byte) & $> 7.9$ & \textbf{7.9996} & 7.9954 & 7.9948 & 7.9949 & & 7.9534 & 7.9484 & 7.9469 & 7.9508 \\
PSNR DC (dB) & $< 10$ & 7.58 & 7.97 & \textbf{7.30} & 8.03 & & 7.73 & 8.04 & \textbf{7.26} & 7.99 \\
SSIM DC & $\approx 0$ & \textbf{0.0031} & 0.0819 & $-$0.0586 & 0.0895 & & \textbf{0.0184} & 0.0855 & $-$0.0715 & 0.0783 \\
TSSIM DC encrypted & $<$ baseline & \textbf{0.0071} & 0.0238 & 0.0217 & 0.0218 & & \textbf{$-$0.0199} & 0.0164 & 0.0082 & 0.0236 \\
\midrule
\multicolumn{2}{l}{Baseline TSSIM DC (original)} & \multicolumn{4}{c}{S1 = 0.0163} & & \multicolumn{4}{c}{S2 = 0.0234} \\
\bottomrule
\end{tabular}%
}
\begin{tablenotes}
\small
\item N/A: không áp dụng cho AES/DES/RC4 (keystream cố định, không có diffusion — xem Chương~\ref{ch:baseline}).
\end{tablenotes}
\end{table}

%% ============================================================
%% Bảng 5: Correctness
%% ============================================================
\begin{table}[ht]
\centering
\caption{Tính đúng đắn — S1 và S2 (kết quả giống nhau)}
\label{tab:correctness}
\begin{tabular}{llcccc}
\toprule
Metric & Ngưỡng & hybrid & aes & des & rc4 \\
\midrule
BER & $= 0$ & 0 & 0 & 0 & 0 \\
PSNR (original$\to$decrypt, dB) & $= \infty$ & $\infty$ & $\infty$ & $\infty$ & $\infty$ \\
Codec compliance (encrypted) & OK & OK$^*$ & OK$^*$ & OK$^*$ & OK$^*$ \\
Codec compliance (decrypted) & PASS & PASS & PASS & PASS & PASS \\
SHA-256 match & PASS & PASS & PASS & PASS & PASS \\
\bottomrule
\end{tabular}
\begin{tablenotes}
\small
\item $^*$ OK + warnings là expected behavior — DC bytes bị mã hoá làm sai H.264 syntax;
FFmpeg dùng error concealment. File decrypted hoàn toàn sạch (no errors).
\end{tablenotes}
\end{table}


\subsection{Tính đúng đắn (Correctness)}

\subsubsection{Kết quả tổng quan}

Tính đúng đắn là điều kiện tiên quyết (prerequisite) của toàn bộ hệ thống:
nếu giải mã không hoàn toàn byte-exact, mọi kết quả bảo mật và hiệu năng đều
vô nghĩa. Kết quả thực nghiệm (Bảng~\ref{tab:correctness}) cho thấy tất cả
8 tổ hợp đều PASS hoàn toàn:

\begin{itemize}
  \item \textbf{BER = 0 tuyệt đối:} Không một bit nào lỗi sau giải mã trong
        tất cả 8 tổ hợp, trên cả S1 (19.037 NALU mã hoá) và S2 (155 NALU).

  \item \textbf{SHA-256 PASS:} Hash của file gốc (219 MB) và file giải mã
        khớp hoàn toàn --- xác suất false positive của SHA-256 là $2^{-256}$,
        thực chất bằng 0. Điều này xác nhận byte-exact ở cấp độ toàn bộ file.

  \item \textbf{PSNR (original$\to$decrypt) = $\infty$:} MSE = 0 (không có
        pixel sai lệch giữa gốc và giải mã sau decode 60 frames).

  \item \textbf{Codec compliance decrypted: PASS:} FFmpeg decode file giải mã
        không phát ra bất kỳ warning nào --- cú pháp H.264 được phục hồi hoàn toàn.
\end{itemize}

\subsubsection{Lý giải tính đúng đắn của thiết kế XOR}

Tính đúng đắn là hệ quả tất yếu của thiết kế: tất cả 4 thuật toán đều dựa trên
XOR thuần túy. Tính chất $a \oplus a = 0$ đảm bảo:

\begin{equation}
\text{Dec}(\text{Enc}(P, K, n), K, n) = P \quad \forall P, K, n
\end{equation}

Miễn là keystream được sinh lại từ cùng $(K, n)$ và hoán vị được đảo ngược chính xác,
giải mã luôn cho kết quả byte-exact. Đây là \textbf{invariant cứng (hard invariant)}
của thiết kế.

\subsubsection{Codec compliance của file encrypted}

Khi file H.264 đã mã hoá DC được decode bằng FFmpeg, kết quả là ``OK + warnings'':
\begin{itemize}
  \item \textbf{OK:} FFmpeg decode được tất cả frames mà không crash.
  \item \textbf{Warnings:} FFmpeg báo syntax errors trong DC region --- hành vi
        dự kiến (expected) vì DC bytes không còn hợp lệ H.264 syntax.
  \item \textbf{Error concealment:} FFmpeg áp dụng cơ chế error concealment
        để hiển thị frame dù DC sai --- đây chính là cơ chế bảo vệ của decoder.
\end{itemize}

File decrypted: \textbf{PASS} (no warnings) --- xác nhận khôi phục hoàn toàn.

\subsection{Hiệu năng (Performance)}

\subsubsection{Chiến lược S1 --- 19.037 NALUs}

Kết quả đầy đủ tại Bảng~\ref{tab:perf_s1}.

\paragraph{Tốc độ thuật toán (DC time/NALU):}

\begin{table}[ht]
\centering
\caption{Xếp hạng tốc độ thuật toán S1 theo DC encrypt time/NALU}
\label{tab:speed_rank_s1}
\begin{tabular}{lrrl}
\toprule
Xếp hạng & Thuật toán & DC encrypt (ns/NALU) & Ghi chú \\
\midrule
1 (nhanh nhất) & RC4    & 769    & KSA đơn giản, không cần AES block \\
2              & AES    & 2.753  & AES-NI hardware; $3.6\times$ chậm hơn RC4 \\
3              & DES    & 23.500 & Software S-box, không có hardware accel \\
4 (chậm nhất) & Hybrid & 29.245 & PLCM (nalu\_index iter) + Arnold + diffusion \\
\bottomrule
\end{tabular}
\end{table}

RC4 nhanh gấp 38 lần hybrid do không phải tính PLCM iteration. Chi phí chính
của hybrid là PLCM: mỗi NALU phải chạy PLCM $n$ lần để đến $x_n$ (tối ưu hoá
bằng state continuity đã giảm xuống $\approx 25$ lần/NALU thay vì $n$ lần).

AES nhanh gấp DES $\sim$8.5 lần nhờ AES-NI. DES và Hybrid có tốc độ gần nhau
vì DES-S-box và PLCM iteration đều là sequential computation nặng.

\paragraph{DC chiếm \% pipeline (Bảng~\ref{tab:dc_percent}):}

Với S1, DC encrypt time chiếm đáng kể tổng pipeline:
\begin{itemize}
  \item \textbf{Hybrid: 36.9\%} pipeline là DC encrypt time. Tức là gần 1/3
        thời gian pipeline là tính toán thuật toán --- bottleneck thực sự.
  \item \textbf{DES: 29.5\%} --- cũng là bottleneck đáng kể.
  \item \textbf{AES: 5.3\%} --- AES-NI nhanh đến mức I/O chiếm 95\%.
  \item \textbf{RC4: 1.35\%} --- gần như không đáng kể.
\end{itemize}

\textbf{Hàm ý cho ứng dụng thực tế (S1):}
Hybrid và DES là hai thuật toán duy nhất mà tốc độ thuật toán \emph{thực sự ảnh hưởng}
đến thời gian pipeline. AES và RC4 nhanh đến mức bottleneck chuyển sang I/O.

\paragraph{Anomaly DES S1 decrypt:}

DES S1 decrypt pipeline time = 8.686 s (vs. encrypt 1.519 s) là bất thường rõ ràng.
Tuy nhiên, DC decrypt time thực tế = 478 ms (nhất quán với encrypt 447 ms).
Nguyên nhân là disk I/O page cache miss: pipeline decrypt phải đọc file encrypted
(file mới, không có trong OS page cache) trong khi pipeline encrypt đọc file gốc
đã được cached. Kiểm tra thêm: chạy decrypt lần 2 (file đã cached) cho pipeline
time $\approx$ 1.5 s --- xác nhận đây là artifact hệ thống.

\subsubsection{Chiến lược S2 --- 155 NALUs}

Kết quả tại Bảng~\ref{tab:perf_s2}.

\paragraph{Hiện tượng pipeline convergence:}

Đây là quan sát quan trọng nhất của chiến lược S2: \textbf{tất cả pipeline times
hội tụ về $0.84$--$1.11$ s}, bất kể thuật toán, mặc dù hybrid chậm hơn RC4 59 lần
(42.962 vs. 724 ns/NALU theo DC).

\textbf{Giải thích định lượng:}

Với 155 NALU và hybrid (42.962 ns/NALU):
\begin{equation}
\text{DC total} = 155 \times 42.962 \text{ ns} = 6.659 \text{ ms}
\end{equation}

Pipeline time hybrid S2 = 0.990 s. Tỷ lệ:
\begin{equation}
\frac{6.659 \text{ ms}}{990 \text{ ms}} \times 100\% = 0.67\%
\end{equation}

DC chiếm 0.67\% tổng pipeline. Phần còn lại (99.33\%) là:
\begin{itemize}
  \item Đọc file 219 MB từ disk.
  \item Scan 19.348 NALU để xác định 155 cần mã hoá.
  \item Ghi file output 219 MB ra disk.
\end{itemize}

Chi phí I/O này giống nhau cho mọi thuật toán, vì pipeline luôn đọc và ghi
toàn bộ 219 MB bất kể số NALU mã hoá. Do đó, dù thuật toán có tốc độ DC
khác nhau 38--59 lần, sự khác biệt đó bị ``chìm'' trong chi phí I/O.

\textbf{Pipeline convergence:} Khi DC $< 1\%$ pipeline, chi phí thuật toán trở nên
không đáng kể. Đây là hiện tượng kinh điển của Amdahl's Law \cite{amdahl1967}:
việc tối ưu hoá một phần nhỏ của hệ thống (DC encryption) không cải thiện
đáng kể tổng thời gian.

\paragraph{Measurement artifact --- DC/NALU S2 cao hơn S1:}

Quan sát thú vị: DC/NALU của hybrid S2 = 42.962 ns, cao hơn S1 = 29.245 ns.
Với cùng thuật toán, tại sao S2 chậm hơn per-NALU?

Giải thích: với 155 NALU (S2), kết quả có variance cao hơn do mẫu nhỏ. Với
19.037 NALU (S1), kết quả là trung bình ổn định hơn (Central Limit Theorem).
Mặt khác, với S2 chỉ 155 NALU, CPU branch predictor và instruction cache
chưa ``warm up'' hoàn toàn so với S1. Đây là measurement artifact, không phải
thuật toán chậm hơn trong S2.

\subsection{Bảo mật (Security)}

\subsubsection{Entropy DC bytes}

\textbf{Kết quả:} Tất cả 8 tổ hợp đạt Entropy $> 7.9$ bits/byte (Bảng~\ref{tab:security}).
Hybrid S1 đạt cao nhất: $7.9996 \approx 8.0$ (lý thuyết tuyệt đối).

\textbf{Giải thích tại sao tất cả đạt $> 7.9$:}
Entropy cao phản ánh tính ngẫu nhiên của keystream, không phải của thuật toán
mã hoá. AES, DES, và RC4 đều có keystream xuất ra gần như ngẫu nhiên (đó là
yêu cầu của PRG --- Pseudo-Random Generator). XOR với keystream ngẫu nhiên
làm DC bytes encrypted phân phối gần đều.

\textbf{Tại sao hybrid cao nhất:}
Per-NALU keystream variation có nghĩa là các NALU khác nhau dùng keystream khác nhau,
làm phân phối DC bytes encrypted đa dạng hơn. Với AES/DES/RC4 fixed keystream,
có thể có partial correlation giữa DC bytes và keystream, làm Entropy giảm nhẹ.

\textbf{S1 vs S2 Entropy:}
Entropy S1 cao hơn S2 vì S1 có nhiều mẫu hơn ($19.037 \times 25 = 475.925$ bytes
vs. $155 \times 25 = 3.875$ bytes). Với mẫu lớn, ước lượng phân phối chính xác hơn.

\subsubsection{PSNR DC}

\textbf{Kết quả:} Tất cả 8 tổ hợp đạt PSNR DC $< 10$ dB.
DES có PSNR DC thấp nhất: S1 = 7.30 dB, S2 = 7.26 dB.

\textbf{Giải thích:} PSNR DC thấp có nghĩa là DC bytes encrypted khác nhiều với
DC bytes gốc (MSE lớn). Tất cả thuật toán đều đạt mức này vì XOR với keystream
ngẫu nhiên đã đủ để làm MSE tăng cao.

DES đạt PSNR DC thấp nhất có thể do keystream DES tương tác đặc biệt tốt với
phân phối DC của video này (vì DC bytes thường là giá trị nhỏ gần 0 sau quantization,
và DES keystream có giá trị phân phối đều hơn RC4 bias).

\textbf{Hàm ý:} PSNR DC $< 10$ dB là điều kiện cần (nhưng không đủ) cho mã hoá tốt.
Tất cả 4 thuật toán đều đạt --- đây không phải điểm phân biệt giữa chúng.

\subsubsection{SSIM DC}

\textbf{Kết quả:}
\begin{itemize}
  \item \textbf{Hybrid:} SSIM DC = $0.0031$ (S1), $0.0184$ (S2) --- gần 0 nhất.
  \item \textbf{AES:} SSIM DC = $0.0819$ (S1), $0.0855$ (S2) --- dương nhẹ.
  \item \textbf{DES:} SSIM DC = $-0.0586$ (S1), $-0.0715$ (S2) --- âm (phản tương quan).
  \item \textbf{RC4:} SSIM DC = $0.0895$ (S1), $0.0783$ (S2) --- dương nhẹ, cao nhất.
\end{itemize}

\textbf{Giải thích sự khác biệt:}
\begin{itemize}
  \item \textbf{Hybrid gần 0 nhất:} Arnold permutation phá vỡ cấu trúc cục bộ
        trước XOR, sau đó chained diffusion propagates sự thay đổi đều đặn.
        Kết quả là DC encrypted không tương quan với DC gốc.

  \item \textbf{DES âm (phản tương quan):} Keystream DES-OFB với IV = 0 có
        đặc tính đặc biệt: tương tác với phân phối DC video tạo ra phản tương quan
        thống kê. Điều này không phải đặc tính thiết kế mà là hiện tượng ngẫu nhiên
        phụ thuộc vào nội dung video.

  \item \textbf{AES/RC4 dương nhẹ:} Keystream cố định XOR với DC bytes giữ nguyên
        một phần cấu trúc tương quan (vì $C[i] = DC[i] \oplus ks_i$ với $ks_i$ hằng số,
        cấu trúc $\text{DC}[i] - \text{DC}[j]$ được bảo toàn nếu $ks_i = ks_j$).
\end{itemize}

\textbf{SSIM DC = 0 là tốt nhất} vì nó có nghĩa là cấu trúc thống kê của DC bytes
bị phá vỡ hoàn toàn --- kẻ tấn công không thể suy luận về DC gốc từ DC encrypted.

\subsubsection{TSSIM DC (Temporal SSIM)}

TSSIM DC baseline (video gốc không mã hoá):
\begin{itemize}
  \item S1 baseline = 0.0163 (tương quan temporal tự nhiên giữa frames liên tiếp)
  \item S2 baseline = 0.0234
\end{itemize}

\textbf{Kết quả:}
\begin{itemize}
  \item \textbf{Hybrid S1:} $0.0071 < 0.0163$ --- Disruption temporal tốt nhất.
        Mã hoá phá vỡ cấu trúc temporal, TSSIM giảm 56\% so với baseline.
  \item \textbf{Hybrid S2:} $-0.0199 < 0.0234$ --- Disruption rất tốt,
        thậm chí tạo phản tương quan temporal.
  \item \textbf{AES/DES/RC4 S1:} $0.0217$--$0.0238 > 0.0163$ --- Không cải thiện
        so với baseline! Mã hoá fixed keystream \emph{tăng} TSSIM so với gốc.
\end{itemize}

\textbf{Giải thích AES/DES/RC4 TSSIM $>$ baseline:}
Với keystream cố định $ks$:
\begin{equation}
\text{TSSIM}(t) = \text{SSIM}(\text{DC}_t \oplus ks, \text{DC}_{t+1} \oplus ks)
\end{equation}

XOR với cùng $ks$ không thay đổi sự tương quan giữa frames liên tiếp về mặt
cấu trúc tương đối --- nếu $\text{DC}_t$ và $\text{DC}_{t+1}$ tương tự nhau
(video ổn định), ciphertext tương ứng cũng tương tự. Thậm chí, XOR fixed có thể
làm tăng SSIM nếu $ks$ ``align'' tốt với sự khác biệt giữa frames.

Kết quả này nhấn mạnh rằng keystream cố định không những không cải thiện mà còn
có thể duy trì/tăng cường cấu trúc temporal --- nguy hiểm cho tấn công phân tích
temporal.

\subsubsection{NPCR và UACI (chỉ hybrid)}

\textbf{Kết quả chi tiết:}

\begin{table}[ht]
\centering
\caption{NPCR và UACI cho thuật toán hybrid}
\label{tab:npcr_detail}
\begin{tabular}{lrrlr}
\toprule
Chiến lược & NPCR (\%) & Ngưỡng (\%) & Đạt? & UACI (\%) \\
\midrule
S1 (19.037 NALUs) & 99.5260 & 99.6 & \textbf{Chưa} (sát mép) & 32.7414 \\
S2 (155 NALUs)    & 99.6903 & 99.6 & \textbf{Đạt} & 32.4751 \\
\bottomrule
\end{tabular}
\end{table}

\textbf{Tại sao S2 đạt ngưỡng nhưng S1 thì không:}

Đây là kết quả đáng ngạc nhiên: S2 (ít NALU hơn) lại đạt NPCR cao hơn S1.

Giải thích: NPCR được trung bình qua tất cả NALU mã hoá. Với S1, bao gồm cả
P/B-frames (Type 1) với phân phối DC bytes đa dạng hơn (motions vectors, residuals
phức tạp hơn I-frame). Một số P/B-frames có DC bytes pattern đặc biệt làm giảm
hiệu ứng chained diffusion tại byte cuối, kéo NPCR trung bình xuống dưới $99.6\%$.

Với S2, chỉ I-frames (Type 5) được mã hoá. I-frames là intra-coded hoàn toàn,
DC bytes có phân phối ``sạch'' hơn, hiệu ứng chained diffusion lan truyền đều đặn
hơn, dẫn đến NPCR cao hơn.

\textbf{UACI chưa đạt $33.46\%$:}

Cả S1 (32.74\%) và S2 (32.48\%) đều chưa đạt ngưỡng lý thuyết $33.46\%$.
Nguyên nhân: fixed points của Arnold Cat Map.

Từ Bảng~\ref{tab:arnold_perm}, byte 0 ($i=0$) và byte 24 ($i=24$) là fixed points
($\pi(0) = 0$, $\pi(24) = 24$). Điều này có nghĩa là $\text{permuted}[0] = \text{DC}[0]$
--- hoán vị Arnold không dịch chuyển byte 0 sang vị trí khác.

Khi flip byte 0 của DC:
\begin{itemize}
  \item $\text{permuted}'[0] = \text{DC}'[0] = \text{DC}[0] \oplus 1$.
  \item Chained diffusion propagates từ vị trí 0 → 24: tất cả $C[0..24]$ thay đổi.
  \item NPCR cao: $\approx 99.6\%$. ✓
  \item Nhưng $|C_1[24] - C_2[24]|$: byte 24 là cộng XOR của tất cả 25 bytes,
        giá trị biến động mạnh, làm UACI trung bình thấp hơn lý thuyết.
\end{itemize}

Fixed points làm một số UACI values lệch, kéo trung bình xuống $\approx 32.5\%$.

\subsection{Thảo luận tổng hợp}

\subsubsection{Trade-off security vs. performance}

\begin{table}[ht]
\centering
\caption{Tổng hợp trade-off security vs. performance cho 4 thuật toán}
\label{tab:tradeoff}
\begin{tabular}{lccccc}
\toprule
Thuật toán & Entropy & SSIM DC & TSSIM $<$ base? & NPCR $>99.6$\% & Tốc độ S1 \\
\midrule
Hybrid & \textbf{7.9996} & \textbf{0.0031} & \textbf{Có} & \textbf{S2: Có} & 29.245 ns \\
AES    & 7.9954 & 0.0819 & Không & N/A & 2.753 ns \\
DES    & 7.9948 & $-0.0586$ & Không & N/A & 23.500 ns \\
RC4    & 7.9949 & 0.0895 & Không & N/A & \textbf{769 ns} \\
\bottomrule
\end{tabular}
\end{table}

\textbf{Hybrid} đạt bảo mật tốt nhất trên mọi chỉ số: Entropy cao nhất, SSIM DC
gần 0 nhất, TSSIM phá vỡ cấu trúc temporal, và là thuật toán duy nhất có NPCR
có ý nghĩa. Cái giá là tốc độ thấp nhất: 29.245 ns/NALU, chậm hơn RC4 38 lần.

Tuy nhiên, với DC encrypt time tuyệt đối S1 = 556 ms trên 19.037 NALU, và
pipeline total = 1.510 s, hybrid vẫn trong ngưỡng chấp nhận được cho ứng dụng
không yêu cầu real-time cứng (live streaming 30 fps cần $< 33$ ms latency).

\textbf{RC4} nhanh nhất (769 ns/NALU) nhưng kém bảo mật nhất: SSIM DC cao (0.09),
không phá vỡ TSSIM, NPCR N/A. RC4 phù hợp khi yêu cầu bảo mật thấp và tốc độ
là ưu tiên hàng đầu.

\textbf{AES} là lựa chọn cân bằng: chuẩn hoá FIPS, nhanh nhờ AES-NI (2.753 ns/NALU),
keyspace $2^{128}$. Thiếu diffusion và per-NALU variation nhưng nếu chỉ cần
Entropy $> 7.9$ và PSNR DC $< 10$ dB thì AES đủ.

\textbf{DES} không khuyến nghị: keyspace $2^{56}$ đã bị phá vỡ, chậm như hybrid
nhưng không có bảo mật của hybrid.

\subsubsection{Ý nghĩa của kết quả NPCR S2}

NPCR S2 = $99.6903\% > 99.6\%$ là kết quả quan trọng nhất của nghiên cứu.
Nó xác nhận rằng cơ chế chained diffusion (Phương trình~\ref{eq:chained_diffusion}
từ Chương 3) hoạt động hiệu quả trên dữ liệu thực: 1 bit thay đổi DC
propagates thay đổi đến $99.69\%$ trong 25 bytes ciphertext.

Kết quả này vượt ngưỡng lý thuyết $99.6\%$ của Wu \cite{wu_npcr} --- chuẩn
thường dùng cho image encryption --- cho thấy thuật toán hybrid có độ nhạy plaintext
tốt ngay cả khi chỉ mã hoá 25 bytes/NALU.

\subsubsection{So sánh với ngưỡng chuẩn của Wu et al.}

Các nghiên cứu mã hoá ảnh thường báo cáo NPCR $>99.6\%$ và UACI $\approx 33.46\%$
trên ảnh $256 \times 256 = 65.536$ bytes. Nghiên cứu này đạt kết quả tương đương
chỉ với \textbf{25 bytes/NALU} --- do cơ chế chained diffusion hoạt động
trên toàn bộ 25 bytes thay vì xử lý từng byte độc lập.

\subsubsection{Hạn chế của kết quả thực nghiệm}

\begin{enumerate}
  \item \textbf{Single dataset:} Tất cả kết quả chỉ trên 1 file video (output.h264).
        Không biết kết quả có tổng quát hoá cho video khác không.

  \item \textbf{Single key:} Chỉ test với \texttt{testkey}. Với key khác, NPCR/UACI
        có thể dao động (phụ thuộc vào $x_0, p$ dẫn xuất từ key).

  \item \textbf{Timing not averaged:} DC time đo từ 1 lần chạy (không lấy trung bình
        nhiều lần). Variance $\pm$5\% có thể ảnh hưởng đến so sánh sát nhau.

  \item \textbf{NPCR S1 chưa đạt ngưỡng:} 99.5260\% $< 99.6\%$ --- chưa đạt chuẩn Wu.
        Cần phân tích thêm với nhiều video và key khác nhau.
\end{enumerate}
